\scanpage{52}

\unsure{général} pour les topos et \unsure{ouverts} d'iceux);
et \unsure{qu'on} a un \unsure{isom.} can.
\inserted{$i_!=(ji')_!=j_!i'_!,\quad i^!=i'^!j^!=i'^!j^*$}
on est ramené à définir \deleted{\unsure{un}} \deleted{\unsure{un} \unsure{foncteur}}
$i'^!$ comme foncteur adjoint de $i'_!$, ce qui
nous ramène au cas où $i$ est une imm.
fermée, et où c'est également un fait général
valable pour les topos et \unsure{ouverts} d'iceux,
\inserted{cf Exp \uline{IV}}

\deleted{Comme} Comme $i_!$ est exact, \deleted{\unsure{on}} il s'ensuit
que son adjoint à droite $i^!$ transforme
\uline{injectifs} en \uline{injectifs}, \deleted{donc} On va
considérer
\[
\mathbb{R}i^!:\quad\mathrm{D}^+(X)\mathbin{\unsure{-|\longrightarrow}}\mathrm{D}^+(Y)
\]
qui, pour $i$ immersion \unsure{ouverte}, n'est
autre que le \uline{foncteur} \uline{restriction}
à $U$ (\deleted{\unsure{aux}}\unsure{sur} cat. dérivées).
\deleted{et pour $i$ immersion \uline{fermée}, n'est autre
que le foncteur $\mathbb{R}\underline{\Gamma}_Y$}

Notons, si
\deleted{\unsure{$\mathbb{R}i$}}
$i$ est une imm. \uline{fermée}
\[
(\mathbb{R}i_{\unsure{*|!}})\deleted{\unsure{i}}(\mathbb{R}i^!)\simeq \mathbb{R}\underline{\Gamma}_Y\deleted{\unsure{)}}
\]
formule \unsure{essentiellement} \unsure{triviale}, \deleted{donc}
si $F\in\mathrm{C}(X)$, on a donc
\[
\underline{H}^{q}(\mathbb{R}i^!(F))\simeq i^*(\underline{H}^{q}_Y(F)).
\]
\mspage{51}
